LLM-based agents increasingly rely on external memory systems to support long-running interactions. However, many real-world memories are not static facts: plans, preferences, project states, and event descriptions often evolve through multiple semantic versions. Existing benchmarks have advanced the evaluation of long-term memory, temporal reasoning, and knowledge update, but they rarely isolate a more specific question: can a memory system retrieve the version memory that matches a queried temporal semantic state from an evolving memory trajectory? To address this gap, we introduce \textbf{TSSMRBench}, a benchmark for temporal semantic state memory retrieval.

TSSMRBench treats stage-correct version retrieval as the target capability of evaluation. We construct a benchmark with 300 scenarios, 9{,}000 memory nodes, and 900 multiple choice questions organized into three task families: \emph{single-state lookup}, \emph{cross-version comparison}, and \emph{temporal ordering}. The tasks differ in how many supporting states they require, enabling evaluation of both retrieval precision and support completeness under shared-pool interference while holding the answer model fixed.

Experiments with BM25, FAISS, Mem0, and Graphiti show that current memory systems perform relatively well on single-state queries but degrade substantially on cross-version comparison and temporal ordering. FAISS is the strongest automatic baseline, while Mem0 and Graphiti expose different trade-offs in version fidelity, support recovery, and latency. These findings show that retrieving evolving memories remains challenging for current memory systems and suggest that future designs should preserve version history, return coherent state-level evidence, and support stage-aware retrieval over semantically similar memories.
